% diversity_reduction_second_order.tex
% Second-order Taylor expansion of the sphere projection under steering.
% Supplements the first-order analysis in diversity_reduction_theory.tex.
%
% Intended to appear as a subsection after the first-order covariance
% derivation, e.g. as Section X.3.1 "Second-order corrections."
%
% === NOTES ===
%
% The main payoffs of the second-order analysis are:
%   1. The mean resultant length is strictly < 1, connecting the Taylor
%      framework to the spherical variance from the geometric analysis.
%   2. The O(1/||m||^3) covariance correction vanishes for symmetric
%      distributions (e.g. Gaussian inputs), so the first-order covariance
%      is actually accurate to O(1/||m||^4) in that case.
%   3. The radial-perpendicular coupling term captures the "magnifying
%      glass" effect near the origin as a perturbative correction.
%
% This section is optional for the main paper and could be deferred to
% an appendix if space is tight. The key results that matter for the
% paper's claims are already in the first-order analysis; this section
% provides rigor and connects the two proof strategies.
%
% === END NOTES ===

\subsection{Second-Order Corrections}
\label{sec:second-order}

The first-order expansion of Section~\ref{sec:covariance} gives the
leading behavior but leaves open the question of how large the
approximation error is.
We now compute the second-order correction, which reveals additional
structure and confirms consistency with the global geometric bounds
of Section~\ref{sec:geometric-bound}.

\paragraph{The Hessian of the sphere projection.}

We require the second derivatives of $g_i(v) = v_i / \|v\|$.
Starting from the Jacobian entry derived in Lemma~\ref{lem:jacobian}:
%
\begin{equation}
  \frac{\partial g_i}{\partial v_j}
  = \frac{\delta_{ij}}{\|v\|} - \frac{v_i v_j}{\|v\|^3},
\end{equation}
%
we differentiate with respect to $v_k$.
The first term contributes $-\delta_{ij}\,v_k / \|v\|^3$ via
$\partial_{v_k}(\|v\|^{-1}) = -v_k / \|v\|^3$.
The second term contributes, by the product and chain rules:
%
\begin{equation}
  -\frac{\delta_{ik}\,v_j + v_i\,\delta_{jk}}{\|v\|^3}
  + \frac{3\,v_i\,v_j\,v_k}{\|v\|^5}.
\end{equation}
%
Combining:
%
\begin{equation}
  \frac{\partial^2 g_i}{\partial v_j \partial v_k}
  = -\frac{\delta_{ij}\,v_k + \delta_{ik}\,v_j + v_i\,\delta_{jk}}{\|v\|^3}
  + \frac{3\,v_i\,v_j\,v_k}{\|v\|^5}.
  \label{eq:hessian-components}
\end{equation}

\paragraph{Contracting with $\delta$.}

The second-order term in the Taylor expansion is
$\frac{1}{2}\sum_{j,k} \frac{\partial^2 g_i}{\partial v_j \partial v_k}\big|_{v=m} \delta_j\,\delta_k$.
Evaluating \eqref{eq:hessian-components} at $v = m$ and summing over
$j$ and $k$, the four groups yield:
%
\begin{align}
  -\frac{\delta_{ij}\,m_k\,\delta_j\,\delta_k}{\|m\|^3}
  &= -\frac{\delta_i\,(\hat{m}^\top\delta)}{\|m\|^2},
  \label{eq:hess-term1}\\[4pt]
  -\frac{\delta_{ik}\,m_j\,\delta_j\,\delta_k}{\|m\|^3}
  &= -\frac{(\hat{m}^\top\delta)\,\delta_i}{\|m\|^2},
  \label{eq:hess-term2}\\[4pt]
  -\frac{m_i\,\delta_{jk}\,\delta_j\,\delta_k}{\|m\|^3}
  &= -\frac{\hat{m}_i\,\|\delta\|^2}{\|m\|^2},
  \label{eq:hess-term3}\\[4pt]
  \frac{3\,m_i\,m_j\,m_k\,\delta_j\,\delta_k}{\|m\|^5}
  &= \frac{3\,\hat{m}_i\,(\hat{m}^\top\delta)^2}{\|m\|^2}.
  \label{eq:hess-term4}
\end{align}
%
Summing and writing in vector form:
%
\begin{equation}
  H_g(m)[\delta,\delta]
  = \frac{1}{\|m\|^2}\Big[
    -2(\hat{m}^\top\delta)\,\delta
    - \|\delta\|^2\,\hat{m}
    + 3(\hat{m}^\top\delta)^2\,\hat{m}
  \Big].
  \label{eq:hessian-vector}
\end{equation}

\paragraph{Decomposition into radial and perpendicular components.}

Let $\alpha = \hat{m}^\top\delta$ denote the radial component of $\delta$
(its projection onto $\hat{m}$), and let
$\delta_\perp = P_{\perp m}\,\delta = \delta - \alpha\,\hat{m}$
denote the perpendicular component.
Then $\|\delta\|^2 = \alpha^2 + \|\delta_\perp\|^2$, and substituting
into \eqref{eq:hessian-vector}:
%
\begin{align}
  H_g(m)[\delta,\delta]
  &= \frac{1}{\|m\|^2}\Big[
    -2\alpha\,(\alpha\,\hat{m} + \delta_\perp)
    - (\alpha^2 + \|\delta_\perp\|^2)\,\hat{m}
    + 3\alpha^2\,\hat{m}
  \Big]
  \notag\\
  &= \frac{1}{\|m\|^2}\Big[
    \underbrace{(-2\alpha^2 - \alpha^2 - \|\delta_\perp\|^2 + 3\alpha^2)}_{=\;-\|\delta_\perp\|^2}\,\hat{m}
    - 2\alpha\,\delta_\perp
  \Big]
  \notag\\
  &= \frac{1}{\|m\|^2}\Big[
    -\|\delta_\perp\|^2\,\hat{m}
    - 2\alpha\,\delta_\perp
  \Big].
  \label{eq:hessian-simplified}
\end{align}
%
The $\alpha^2$ terms in the $\hat{m}$ direction cancel exactly.

\paragraph{Full second-order expansion.}

Combining the first-order term from \eqref{eq:taylor} with
$\frac{1}{2}H_g(m)[\delta,\delta]$:
%
\begin{equation}
  z \approx \hat{m}
  + \frac{1}{\|m\|}\,\delta_\perp
  + \frac{1}{2\|m\|^2}\Big[
    -\|\delta_\perp\|^2\,\hat{m}
    - 2\alpha\,\delta_\perp
  \Big]
  + O\!\left(\frac{\|\delta\|^3}{\|m\|^3}\right).
  \label{eq:taylor-second}
\end{equation}
%
Regrouping by direction:
%
\begin{equation}
  \boxed{
  z \approx
  \left(1 - \frac{\|\delta_\perp\|^2}{2\|m\|^2}\right)\hat{m}
  \;+\;
  \left(\frac{1}{\|m\|} - \frac{\alpha}{\|m\|^2}\right)\delta_\perp
  \;+\;
  O\!\left(\frac{\|\delta\|^3}{\|m\|^3}\right).
  }
  \label{eq:taylor-second-grouped}
\end{equation}
%
Two effects appear at second order:
%
\begin{enumerate}
  \item \textbf{Radial pullback.}
    The coefficient of $\hat{m}$ decreases from $1$ by
    $\|\delta_\perp\|^2 / (2\|m\|^2)$.
    This is the sphere curvature: spreading perpendicular to the pole
    pulls the projection inward, reducing the mean resultant length.
    This is precisely the mechanism underlying the spherical variance
    bound of Corollary~\ref{cor:spherical-variance}.

  \item \textbf{Radial-perpendicular coupling.}
    The coefficient of $\delta_\perp$ acquires a correction
    $-\alpha / \|m\|^2$ that depends on the radial fluctuation
    $\alpha = \hat{m}^\top\delta$.
    Points with $\alpha > 0$ (further from the origin) receive
    \emph{less} angular spread; points with $\alpha < 0$ (closer to
    the origin) receive \emph{more}.
    This is the perturbative manifestation of the magnifying-glass
    effect discussed in Section~\ref{sec:anti-pole}: the normalization
    map amplifies angular variation near the origin.
\end{enumerate}

\paragraph{Second-order mean.}

Taking expectations using $\mathbb{E}[\delta] = 0$ (and therefore
$\mathbb{E}[\alpha] = 0$, $\mathbb{E}[\delta_\perp] = 0$):
%
\begin{equation}
  \mathbb{E}[z]
  \approx \hat{m}
  - \frac{\mathbb{E}[\|\delta_\perp\|^2]}{2\|m\|^2}\,\hat{m}
  - \frac{\mathbb{E}[\alpha\,\delta_\perp]}{\|m\|^2}.
  \label{eq:mean-second}
\end{equation}
%
For the first correction:
$\mathbb{E}[\|\delta_\perp\|^2]
= \mathrm{tr}(P_{\perp m}\,\Sigma_x\,P_{\perp m})
= \mathrm{tr}(\Sigma_x) - \hat{m}^\top\Sigma_x\,\hat{m}$,
which is the perpendicular variance from
Proposition~\ref{prop:diversity-reduction}.
For the second correction:
$\mathbb{E}[\alpha\,\delta_\perp]
= \mathbb{E}[(\hat{m}^\top\delta)\,P_{\perp m}\,\delta]
= P_{\perp m}\,\Sigma_x\,\hat{m}$,
which is the cross-covariance between the radial and perpendicular
components of $\delta$.
This vanishes when $\hat{m}$ is an eigenvector of $\Sigma_x$
or when $\Sigma_x$ is isotropic.

The mean resultant length is therefore:
%
\begin{equation}
  \|\mathbb{E}[z]\|
  \approx 1
  - \frac{\mathrm{tr}(\Sigma_x) - \hat{m}^\top\Sigma_x\,\hat{m}}{2\|m\|^2}
  + O\!\left(\frac{\|\Sigma_x\|^2}{\|m\|^4}\right),
  \label{eq:resultant-length}
\end{equation}
%
where the cross-covariance term contributes only at $O(1/\|m\|^4)$
to the norm (since it is perpendicular to the leading $\hat{m}$ term).
The spherical variance is thus:
%
\begin{equation}
  V \coloneqq 1 - \|\mathbb{E}[z]\|
  \approx \frac{\mathrm{tr}(\Sigma_x) - \hat{m}^\top\Sigma_x\,\hat{m}}{2\|m\|^2},
  \label{eq:spherical-variance-taylor}
\end{equation}
%
which scales as $O(1/\|m\|^2) = O(1/\|s\|^2)$ in the strong-steering
regime, consistent with the $O(1/\|s\|)$ chordal bound from
Corollary~\ref{cor:spherical-variance} (since spherical variance scales
as the square of the chordal spread for concentrated distributions).

\paragraph{Second-order covariance.}

The fluctuation around the corrected mean is, to second order:
%
\begin{equation}
  \varepsilon = z - \mathbb{E}[z]
  = \frac{1}{\|m\|}\,\delta_\perp
  - \frac{1}{\|m\|^2}\Big[\alpha\,\delta_\perp
    - \mathbb{E}[\alpha\,\delta_\perp]\Big]
  - \frac{1}{2\|m\|^2}\Big[\|\delta_\perp\|^2
    - \mathbb{E}[\|\delta_\perp\|^2]\Big]\,\hat{m}
  + O\!\left(\frac{\|\delta\|^3}{\|m\|^3}\right).
  \label{eq:fluctuation-second}
\end{equation}
%
The covariance $\Sigma_z = \mathbb{E}[\varepsilon\,\varepsilon^\top]$
has three tiers:
%
\begin{enumerate}
  \item \textbf{Leading term}, $O(1/\|m\|^2)$:
    $\frac{1}{\|m\|^2}\,P_{\perp m}\,\Sigma_x\,P_{\perp m}$,
    exactly as in Proposition~\ref{prop:diversity-reduction}.

  \item \textbf{Cross terms}, $O(1/\|m\|^3)$:
    these involve third central moments of $\delta$, specifically
    $\mathbb{E}[\delta_\perp \cdot \alpha\,\delta_\perp^\top]$
    and its transpose.
    \emph{These vanish for any distribution symmetric about $\mu$.}
    In particular, for Gaussian inputs the first-order covariance is
    accurate to $O(1/\|m\|^4)$, not merely $O(1/\|m\|^3)$.

  \item \textbf{Quadratic terms}, $O(1/\|m\|^4)$:
    these involve fourth central moments such as
    $\mathbb{E}[\alpha^2\,\delta_\perp\,\delta_\perp^\top]$
    and $\mathbb{E}[\|\delta_\perp\|^4]$.
    For Gaussian inputs these can be expressed via Isserlis' theorem
    in terms of $\Sigma_x$, but the resulting expressions are
    cumbersome and we omit them.
\end{enumerate}

\paragraph{Summary.}
The second-order analysis confirms and sharpens the first-order results:
%
\begin{center}
\begin{tabular}{lcc}
  \toprule
  \textbf{Quantity} & \textbf{First-order} & \textbf{Second-order correction} \\
  \midrule
  Mean direction & $\hat{m}$ & unchanged \\[4pt]
  Mean resultant length & $1$ &
    $1 - \dfrac{\mathrm{tr}(\Sigma_x) - \hat{m}^\top\Sigma_x\hat{m}}{2\|m\|^2}$ \\[10pt]
  Spherical variance & $0$ (trivially) &
    $\dfrac{\mathrm{tr}(\Sigma_x) - \hat{m}^\top\Sigma_x\hat{m}}{2\|m\|^2}$ \\[10pt]
  Covariance &
    $\dfrac{P_{\perp m}\,\Sigma_x\,P_{\perp m}}{\|m\|^2}$ &
    same $+ O(1/\|m\|^4)$ for symmetric $\delta$ \\[6pt]
  \bottomrule
\end{tabular}
\end{center}
%
\noindent
The key takeaway is that for symmetric input distributions (including
Gaussian), the first-order covariance formula from
Proposition~\ref{prop:diversity-reduction} is more accurate than its
derivation might suggest: the next correction is $O(1/\|m\|^4)$, not
$O(1/\|m\|^3)$.
For asymmetric distributions, third-moment effects enter at
$O(1/\|m\|^3)$ and can either increase or decrease the covariance
depending on the skewness structure.
